\section{Scales, energies, and limiting arguments}
\label{sec:framework}

The arguments below compare experiments whose parameters and probability
spaces vary with scale.  This section records the meaning of those
comparisons.  It does not assert that a geometrically defined
subexperiment is admissible: its support, regularity, and operator bounds
must be checked in the section where it is constructed.

\subsection{Fixed setups and exponent bounds}

We use dyadic scales $s=2^{-n}$ tending to zero.  A depth $a$ represents
the scale $s^a$.  Replacing it by the nearest dyadic scale changes a
positive quantity by a bounded factor and hence changes its normalized
logarithm by $O(1/\log(1/s))$.  Width exponents and growth exponents have
opposite sign conventions:
\[
 \operatorname{width\ exponent}(r)=\frac{\log r}{\log s},
 \qquad
 \operatorname{growth\ exponent}(B)=\frac{\log B}{\log(1/s)}.
\]
The growth exponent of a zero quantity is $-\infty$.

\begin{definition}[Fixed polynomial setup]\label{def:fw-setup}
A fixed polynomial setup specifies all structural constants, the finite
upper bounds on powers of $s^{-1}$ occurring in its hypotheses, and a
finite bound on the number of operations along each branch.  These data
are fixed before $s\to0$.  A sequence in this setup may have parameters,
operators, measures, and Hilbert spaces depending on $s$, provided the
specified bounds hold uniformly along the sequence.

An error $r_s$ is negligible relative to a positive reference quantity
$a_s$ if, for every $N>0$,
\[
 |r_s|\le C_N s^N a_s
\]
for all sufficiently small scales, with $C_N$ independent of scale.
For a family of choices at each scale, uniform negligibility means that
the same constants and scale thresholds work for every choice.  The
reference quantity is always specified; it may be an input energy or
an operator norm.  Polynomially many errors can be added without losing
negligibility when their reference quantities and polynomial bounds
are controlled uniformly.
\end{definition}

For a nonnegative functional $\mathcal B_s$, define its upper exponent
to be
\begin{equation}\label{eq:fw-upper-exponent}
 \Gamma=\sup_{\mathscr S}\ \sup_{(s_k,\mathcal E_k)}
 \limsup_{k\to\infty}
 \frac{\log\mathcal B_{s_k}(\mathcal E_k)}{\log(1/s_k)}.
\end{equation}
Here $\mathscr S$ ranges over fixed polynomial setups, and the inner
supremum ranges over admissible sequences in that setup with $s_k\to0$.
The particular model specifies which setups are allowed.  The finiteness
of $\Gamma$ must be proved for that model.

\begin{lemma}[Uniform bound within a fixed setup]
\label{lem:fw-uniform}
Suppose $\Gamma\in\R$ in \eqref{eq:fw-upper-exponent}.  Fix a family
$\mathcal A_s$ such that every sequence of choices
$\mathcal E_k\in\mathcal A_{s_k}$, $s_k\to0$, is admissible in one fixed
polynomial setup.  In particular, all negligible tail hypotheses are
uniform over this family.  For every $\eps>0$ there is $s_0>0$ such that
\begin{equation}\label{eq:fw-uniform-bound}
 \mathcal B_s(\mathcal E)\le s^{-\Gamma-\eps}
 \qquad(0<s<s_0,\ \mathcal E\in\mathcal A_s).
\end{equation}
If the functional is also uniformly bounded on this family at each of
the finitely many larger dyadic scales under consideration, then
\eqref{eq:fw-uniform-bound} holds on the entire scale range after
inserting a constant $C_{\eps,\mathscr S}$.
\end{lemma}

\begin{proof}
Failure of the first assertion provides $s_k\to0$ and
$\mathcal E_k\in\mathcal A_{s_k}$ with
$\mathcal B_{s_k}(\mathcal E_k)>s_k^{-\Gamma-\eps}$.
The resulting sequence is admissible in the fixed setup and has upper
exponent at least $\Gamma+\eps$, contradicting
\eqref{eq:fw-upper-exponent}.  For the second assertion, take the
maximum of the finitely many required constants at the remaining
scales and $1$.
\end{proof}

The constant in Lemma~\ref{lem:fw-uniform} may depend on every bound
fixed in Definition~\ref{def:fw-setup}.  It is not uniform when those
bounds grow along an outer approximation.  Conversely, a failure of an
estimate with a fixed positive exponent slack can be disproved within
one fixed setup, even when that setup was selected late in an outer
argument.

\begin{remark}[Small normalized pieces]\label{rem:fw-small-pieces}
Changing the reference energy can change a tail hypothesis.  For
example, an error negligible relative to $e$ is negligible relative to
$e'$ whenever $e'/e\ge s^C$ for some fixed $C$, but this implication
does not hold for arbitrary $e'>0$.  Applications to a collection of
subexperiments must therefore either verify the hypotheses relative
to each new energy, or use a direct estimate on pieces with
superpolynomially small energy.  No uniformity assertion in
Lemma~\ref{lem:fw-uniform} supplies this missing verification.
\end{remark}

\subsection{Energy and finite extraction}

Energy means a squared Hilbert norm.  Input energy at a time interval
means the energy before the evolutions and masks being estimated.
When a time interval is subdivided, its input is copied into its
descendants.  Consequently the sum of input energies over distinct
starting intervals need not be bounded by the original energy.
This sum will always be distinguished from the energy available to
one interval.

\begin{lemma}[Disjoint row restrictions and copies]
\label{lem:fw-energy}
Let $V:\mathcal H\to\mathcal K$ be an isometry and let $P_1,\ldots,P_N$
be mutually orthogonal projections on $\mathcal K$.  For
$g_j=V^*P_jVg$,
\begin{equation}\label{eq:fw-row-energy}
 \sum_{j=1}^N\norm{g_j}_{\mathcal H}^2
 \le\sum_{j=1}^N\norm{P_jVg}_{\mathcal K}^2
 \le\norm g_{\mathcal H}^2.
\end{equation}
If $\sum_jP_jVg=Vg$, then $\sum_jg_j=g$.
By contrast, for $N$ contractions $T_j$ applied to separate copies of
$g$, the general bound is only
$\sum_j\norm{T_jg}^2\le N\norm g^2$.
\end{lemma}

\begin{proof}
The adjoint $V^*$ has norm one.  Apply this fact to each $P_jVg$ and
then use orthogonality.  The reconstruction assertion follows from
$V^*V=\Id$.  The assertion about copies follows by applying
$\norm{T_jg}\le\norm g$ separately to each $j$.
\end{proof}

The projections in Lemma~\ref{lem:fw-energy} may be multiplication by
disjoint measurable row sets in a direct integral.  Their synthesized
images need not be orthogonal.  For the weighted sums used here, this
is accommodated by a norm inequality, not by an orthogonality claim.

\begin{lemma}[Extraction from finitely many classes]
\label{lem:fw-finite-extraction}
Let $q\ge1$, and suppose that a nonnegative functional $F$ satisfies
\[
 F(g_1+\cdots+g_N)^{1/q}
 \le\sum_{j=1}^N F(g_j)^{1/q}.
\]
Then some $j$ satisfies $F(g_j)\ge N^{-q}F(\sum_i g_i)$.
Consequently, if $N=N_s$ obeys
$\log N_s/\log(1/s)\to0$, passing to a maximizing class costs zero
growth exponent.  The same conclusion holds for a fixed finite
number of successive such extractions.  To use the conclusion for
a model exponent, each selected class must itself be admissible in
a fixed setup.

In particular, for linear maps $A_\lambda$ into Hilbert spaces and
positive weights $w_\lambda$,
\begin{equation}\label{eq:fw-cube-seminorm}
 F(g)=\sum_\lambda w_\lambda^{-1/2}\norm{A_\lambda g}^3
\end{equation}
has the required inequality with $q=3$.
\end{lemma}

\begin{proof}
The right side of the assumed inequality is at most
$N\max_jF(g_j)^{1/q}$.  Raising to the $q$th power proves the
extraction estimate.  Its logarithmic loss is
$q\log N_s/\log(1/s)$; finitely many such losses still tend to zero.
For \eqref{eq:fw-cube-seminorm}, $F(g)^{1/3}$ is the norm of
$(w_\lambda^{-1/6}A_\lambda g)_\lambda$ in the corresponding
$\ell^3$ sum of Hilbert spaces, and Minkowski's inequality applies.
For a countable label set, apply the inequality first to finite
subsets and then use monotone convergence.
\end{proof}

A separate elementary extraction will be used when energies are
already summable: if $0<E<\infty$, $A>0$, $a_j\ge0$, $e_j\ge0$,
$\sum_j e_j\le E$, and $\sum_j a_j\ge A$, then
$\sup_{j:e_j>0}a_j/e_j\ge A/E$, provided $a_j=0$ whenever $e_j=0$.
Indeed the contrary strict upper bound, multiplied by $e_j$ and
summed, contradicts the two hypotheses.  Any additional count or
regularity constraints on the chosen class require their own
discard estimates.

\subsection{Closure of exponent data}

For a finite list of positive quantities, exponent data are the limits
of their normalized logarithms along admissible sequences, when these
limits exist and are finite.  We may also take the closure of such data
in a fixed finite-dimensional Euclidean space.  This outer closure
permits a different fixed polynomial setup for each approximating
sequence.  It does not turn these sequences into one admissible
sequence with unbounded setup parameters.

\begin{lemma}[Extrema of bounded exponent data]
\label{lem:fw-closure}
Let $K\subset\R^{m+1}$ be a nonempty compact set of feasible limiting
data $(g,x_1,\ldots,x_m)$.  Set $G=\max_K g$, minimize $x_1$ on
$K\cap\{g=G\}$, and then successively minimize $x_j$ on the slice
where all preceding extrema are attained.  Every slice is nonempty
and compact, and every stated minimum is attained.

The attainment assertions, and nonemptiness and closedness of the
successive slices, remain valid for a possibly noncompact closed
set $K$ under the following hypotheses: $g$ has a finite supremum
and a bounded maximizing sequence in $K$, and each successive
coordinate has a finite infimum and a bounded minimizing sequence
within the preceding extremal slice.  This alternative does not
assert compactness of the slices.  A new coordinate may be added
to the tuple only after the required boundedness and feasibility
have been proved for the enlarged tuple.
\end{lemma}

\begin{proof}
A continuous coordinate function attains its extremum on a nonempty
compact set.  The set where that extremum is attained is closed in
the same compact set and is nonempty.  Induction proves the first
assertion.  For the alternative, pass first to a convergent
subsequence of the bounded maximizing sequence.  Closedness and
continuity attain the supremum of $g$.  At each subsequent step,
pass to a convergent subsequence of the stipulated bounded
minimizing sequence.  Closedness preserves feasibility, and
continuity of the coordinates preserves the preceding equalities
and attains the next infimum.  Each resulting slice is a nonempty
intersection of closed sets.
\end{proof}

\begin{remark}\label{rem:fw-extra-coordinate}
Boundedness of a new coordinate cannot be inferred from compactness
of previously chosen coordinates.  For example, the closure of the
projection of
$\{(p,c):c\ge1,\ p=1/c\}$ onto the $p$ axis contains $p=0$, although
the closure of the full set has no finite point with $p=0$.
Thus a geometric bound on the curvature cost at the already chosen
growth and mass exponents is an essential part of a curvature
minimization argument.
\end{remark}

\begin{lemma}[Strict improvement in a closure]
\label{lem:fw-strict-improvement}
Let $K\subset\R^{k+1}$ be a closed set of feasible data.  Suppose
that $c_0$ is the minimum of the last coordinate $c$ on the nonempty
slice $K\cap\{z=z_0\}$.  There cannot be data $(z_n,c_n)\in K$ with
$z_n\to z_0$ and $\limsup_n c_n<c_0$, provided $(c_n)$ is bounded
below and has a bounded subsequence.  In particular, a construction
contradicts minimality when its preceding exponents converge to the
specified extrema and its cost improves by a fixed positive amount
in the units of the constructed experiment.
\end{lemma}

\begin{proof}
Pass to a convergent subsequence of $(c_n)$.  Its limit $c$ is strictly
less than $c_0$.  Closedness places $(z_0,c)$ in the minimizing slice,
which is a contradiction.
\end{proof}

These lemmas assert neither an exactly extremizing input nor an
infinite chain of operators.  A construction at limiting exponent
data must be performed on approximating finite experiments.  The
estimates must then imply the asserted bounds on a convergent
subsequence of the finite tuples actually used.

\subsection{Ordered limits and changes of scale}

At a depth gap $h>0$, the new basic scale is $\delta=s^h$.  A factor
$s^{-\alpha}$ has growth exponent $\alpha/h$ in these new units.
Consequently a loss tending to zero in the original units is harmless
on a fixed gap, but it is harmless on shrinking gaps only after it has
been made $o(h)$.

All uses of mesh refinement or tangent gaps have the following order.
First fix the positive gap, the finite list of cuts and observations,
and all cutoff depths needed for that comparison.  In each fixed
polynomial setup take the scale limit.  Next make the preceding
approximation errors and mesh errors as small as required relative
to that fixed gap.  Only then shrink the gap or enlarge the finite
list.  When further tangent layers are used, their positive lengths
and required accuracies are fixed before the preceding tangent limit
is taken.  Each application of a model exponent is to an admissible
sequence at a fixed stage of these outer choices.

\begin{lemma}[Pullback of a strict exponent contradiction]
\label{lem:fw-pullback}
Let $\Gamma\in\R$ be an upper exponent as in
\eqref{eq:fw-upper-exponent}.  Suppose that, for each integer $m$, a
construction produces an admissible sequence in a fixed setup
$\mathscr S_m$ with upper exponent at least
$\Gamma+\tau-r_m$, where $\tau>0$ is fixed and $r_m\to0$.
Then such a construction is impossible.  The setup need not be
uniform in $m$.
\end{lemma}

\begin{proof}
Choose one fixed $m$ with $r_m<\tau/2$.  The sequence in
$\mathscr S_m$ then contradicts the definition of $\Gamma$.
\end{proof}

Having one favorable scale for each $m$ would not satisfy the
hypothesis of Lemma~\ref{lem:fw-pullback}.  For example, the quantities
$B_m(2^{-n})$, equal to $2^n$ when $n=m$ and equal to $1$ otherwise,
have upper exponent zero for every fixed $m$ but exponent one along
the diagonal $m=n$.  We will therefore use a diagonal of scales for
probability comparisons only after the necessary model estimates
with strict slack have been established in their fixed setups.

\subsection{Elementary probability comparisons}

All random variables below are evaluated under the probability law
specified at that point of the argument.  A conditional probability
may condition on a continuous variable or a sigma-algebra. We use
natural logarithms. Under a law $P$, for a finite-valued name $X$,
write
\[
 H_P(X\mid Y)=-\E_P\log P(X\mid Y).
\]
Relative entropy is
$\operatorname{KL}(P\Vert Q)=\int\log(dP/dQ)\,dP$ when $P\ll Q$,
and is $+\infty$ otherwise. Conditional mutual information is
\[
 I_P(X;V\mid Y)
 =\E_{P_Y}\operatorname{KL}
   \bigl(P_{X,V\mid Y}\Vert P_{X\mid Y}\otimes P_{V\mid Y}\bigr).
\]
This definition also permits continuous variables; when $X$ is finite,
it equals $H_P(X\mid Y)-H_P(X\mid V,Y)$. For $0<\rho<1$, the
scale-normalized quantities are $H_\rho=H_P/\log(1/\rho)$ and
$I_\rho=I_P/\log(1/\rho)$, with the law indicated when it changes.
For a finite-valued name $X$ and conditioning data $Y$, write
\begin{equation}\label{eq:fw-score}
 i_s(X\mid Y)
   =-\frac{\log\Prob(X\mid Y)}{\log(1/s)},
\end{equation}
where the conditional probability is evaluated at the sampled value
of $X$. Values on null events are immaterial. Thus
$H_s(X\mid Y)=\E_P i_s(X\mid Y)$: the score is random, whereas
entropy and mutual information are averaged quantities.

\begin{lemma}[Capacity tails]\label{lem:fw-capacity}
Suppose that, given $Y$, the variable $X$ has at most $N$ possible
values almost surely.  For every $u\ge0$,
\begin{equation}\label{eq:fw-capacity}
 \Prob\bigl\{\Prob(X\mid Y)<e^{-u}/N\bigr\}\le e^{-u}.
\end{equation}
If $N\le s^{-C}$ for a fixed $C$, then
$\Prob\{i_s(X\mid Y)>C+t\}\le s^t$ for $t\ge0$.
These normalized scores are uniformly integrable as $s\to0$.
In particular, convergence in probability to a finite constant
implies convergence of their expectations to that constant.
\end{lemma}

\begin{proof}
For each conditioning value, sum the conditional probabilities of
the at most $N$ values below $e^{-u}/N$, and then integrate.  This
proves \eqref{eq:fw-capacity} and its normalized version.  For $K>C$,
the tail integral identity gives
\[
 \E\bigl[i_s(X\mid Y)\mathbf 1_{\{i_s(X\mid Y)>K\}}\bigr]
 \le s^{K-C}\left(K+\frac1{\log(1/s)}\right).
\]
For $s\le1/2$ the right side tends to zero uniformly as $K\to\infty$.
This is uniform integrability; truncation at $K$ proves the final
assertion.
\end{proof}

The same estimate applies to a relative name with at most $N$
refinements after its parent has been included in the conditioning.
It does not require a bounded number of values for the conditioning
data.  When absolute names have unbounded capacity exponents along
an outer limit, the lemma supplies uniform integrability only for
the relative names whose capacities have been bounded.

\begin{lemma}[Stability under restriction]
\label{lem:fw-restriction}
Let $Q=P(\,\cdot\mid E)$, where $P(E)=q>0$, and let $X$ be a finite
name with arbitrary conditioning data $Y$.  If $q\ge s^\alpha$,
then, outside a set of $Q$-probability at most $2s^\tau$,
\begin{equation}\label{eq:fw-restriction}
 \abs{i_s^Q(X\mid Y)-i_s^P(X\mid Y)}\le\alpha+\tau
 \qquad(\tau>0).
\end{equation}
Consequently, restrictions with $q=s^{o(1)}$ preserve deterministic
limiting scores, provided their previous exceptional probabilities
are negligible relative to $q$.  On a gap with scale $s^h$, the
corresponding condition is $\log(1/q)=o(h\log(1/s))$ and all errors
are divided by $h$.
\end{lemma}

\begin{proof}
For the marginal of any variable $W$, the likelihood ratio
$L_W=dQ_W/dP_W$ is $P(E\mid W)/q$, hence is at most $1/q$.
Moreover
$Q\{L_W<s^\tau\}\le s^\tau$ by integrating $L_W$ against $P_W$.
Apply these two facts to $W=(X,Y)$ and $W=Y$.  Except on the union
of the two exceptional sets, both ratios belong to
$[s^\tau,s^{-\alpha}]$.  The conditional likelihood ratio is
$L_{(X,Y)}/L_Y$, so it belongs to
$[s^{\alpha+\tau},s^{-\alpha-\tau}]$.
Taking logarithms proves \eqref{eq:fw-restriction}.
An old exceptional event of $P$-probability $r$ has $Q$-probability
at most $r/q$, proving the stated qualification and the consequences.
\end{proof}

\begin{lemma}[Negligible information and positive restrictions]
\label{lem:fw-information}
Let $U,V$ be finite names, with arbitrary conditioning data $Y$.
Put $L=\log(1/s)$ and
\[
 j=\log\frac{P(U\mid V,Y)}{P(U\mid Y)}.
\]
Then $P\{j<-r\}\le e^{-r}$ for $r\ge0$, and
$\E j=I_P(U;V\mid Y)$.  If $I_P(U;V\mid Y)=o(L)$, then
$j/L\to0$ in probability.  If $Q=P(\,\cdot\mid E)$ and
$P(E)\ge q_0>0$, then
\begin{equation}\label{eq:fw-information-restriction}
 I_Q(U;V\mid Y)
 \le\frac{I_P(U;V\mid Y)+\log2}{q_0}.
\end{equation}
In particular such restrictions preserve $o(L)$ information budgets.
No such conclusion is asserted for arbitrary events of probability
$s^{o(1)}$.
\end{lemma}

\begin{proof}
Condition on $Y$ and compare the joint law of $(U,V)$ with the
product of its conditional marginals.  At actual samples their
likelihood ratio is $e^j$, so integrating over the set where it is
less than $e^{-r}$ proves the tail bound.  Therefore
$\E(j_-)\le1$, while the definition of conditional mutual
information gives $\E j=I_P(U;V\mid Y)$.  It follows that
$\E(j_+)\le I_P(U;V\mid Y)+1$.  Markov's inequality and the negative
tail bound prove convergence in probability after division by $L$.

Let $B=\mathbf 1_E$.  The chain rule gives
\[
 I_P(U;V\mid Y,B)
 \le I_P(U;V,B\mid Y)
 \le I_P(U;V\mid Y)+H_P(B\mid V,Y)
 \le I_P(U;V\mid Y)+\log2.
\]
The left side is the average, over the two values of $B$, of the
corresponding conditional mutual informations.  Its contribution
from $B=1$ is $P(E)I_Q(U;V\mid Y)$, proving
\eqref{eq:fw-information-restriction}.
\end{proof}

The elementary comparisons in this subsection refer to actual laws
and finite lists of names.  Later score limits may be obtained from
joint weak subsequences of those finite lists; no convergence of
conditional sampling kernels is required or asserted.  Extending
the limits to a dense set of depth queries requires the specific
nesting and list bounds proved for those names.  In particular,
these probability lemmas neither couple input and output packet
marginals through a quantum mask nor transfer a profile between
different measuring laws without a separate argument.
